\section*{Chapter \ref{ch:m2:getting-access-fx} --- Getting Access: FX}

\begin{solution}{exo:m2:getting-access-fx:1}
USD~40 million, the sum of the longs. Every trade is long one currency and short
another of the same dollar value, so, counting the dollar as a currency, the net longs
and net shorts always add up to the same amount: here $25 + 15 = 40$.
\end{solution}

\begin{solution}{exo:m2:getting-access-fx:2}
That the client is authorised to trade on the prime broker's behalf with that dealer,
which trade types, currencies and maximum tenors are allowed, through which offices,
and the net open position and \omterm{def:m2:getting-access-fx:limits}{settlement limits} beyond which the prime broker will not
accept the trades.
\end{solution}

\begin{solution}{exo:m2:getting-access-fx:3}
The 2015 franc shock showed the risk in the business when its profitability was
already low and post-crisis rules raised its cost; prime brokers kept large,
profitable clients and dropped retail aggregators, smaller hedge funds and some
high-frequency firms, many of which went to \omterm{def:m2:getting-access-fx:pop}{prime-of-prime} providers.
\end{solution}

\begin{solution}{exo:m2:getting-access-fx:4}
At one prime broker the trades net: the fund receives EUR~40 million, worth USD~44
million, and pays dollars, so the settlement amount is 44 million. At two brokers they
do not: one must receive EUR worth 110 million, the other USD~66 million, 176 million in
all.
\end{solution}

\begin{solution}{exo:m2:getting-access-fx:5}
Because it lowers the prime broker's exposure: refusing it would keep the position above
the limit. The Master Agreement's test is whether a trade causes a limit to be exceeded
or further exceeded.
\end{solution}

\begin{solution}{exo:m2:getting-access-fx:6}
The share of its quotes that are hit and filled, its \omterm{def:m2:fx-spot-microstructure:lastlook}{reject rate} and hold times under
\omterm{def:m2:fx-spot-microstructure:lastlook}{last look}, the mark-outs of trades done against it, the stability of its prices and
sizes, and whether its practices match its disclosure cover sheet.
\end{solution}

\begin{solution}{exo:m2:getting-access-fx:7}
C receives USD~485 million of flow instead of 435, 50 million more, and the day's
position ends with 279 million at C instead of 229, 12 million at A and 120 million at
B; no order is rejected.
\end{solution}

\begin{solution}{exo:m2:getting-access-fx:8}
Liquidity is not the constraint; credit is. One prime broker's limits cap the fund's
position and settlement, it can cut them or exit the business in a shock, as after
2015, and it sees all the fund's flow and can price accordingly. Several prime brokers
give capacity, a fallback and competition, at the cost of more documentation and of
positions spread across them.
\end{solution}

\begin{solution}{pb:m2:getting-access-fx:1}
\textbf{1.} EURUSD and USDJPY: A, B or C; GBPUSD: B or C; USDMXN: C only.
\textbf{2.} USD~55 million both times: the second trade shrinks the dollar short and adds
a yen short.
\textbf{3.} Not for the same value date: settlement would reach 605 million, above 600;
yes for the next day.
\textbf{4.} Only those that reduce the position; any order that adds to it is rejected
until the position is below 50 million.
\textbf{5.} Because the prime broker's exposure grows with the tenor: longer trades are
harder to price, hedge and close out if the client fails.
\textbf{6.} Pesos to C (50); sterling to B (150); the majors to A up to its 500, then B up
to its 600 (450 more), and the last 50 to C: A 500, B 600, C 100.
\textbf{7.} Pesos to C (30); the rest to B (150), A (100) and C (20): A 100, B 150, C 50.
\textbf{8.} Fees USD~5\,300 and overnight charges USD~8\,500: 13\,800 a day, USD~3.45
million a year.
\textbf{9.} USD~24\,600 a day: 9\,600 of fees and 15\,000 of charges.
\textbf{10.} The search finds the same minimum, 13\,800: with one price per broker and
no interaction between the limits, filling the most constrained needs first from the
cheapest broker is optimal.
\textbf{11.} At order 107.
\textbf{12.} None.
\textbf{13.} USD~51 million at A, 150 at B (full) and 229 at C: the router sends each
order to the cheapest fee, and once A and B are full the position builds at C, which
charges most for holding it.
\textbf{14.} Because the gate nets receipts and payments per value date: buys and sells
in the same currency offset, so the settlement amount is far below the gross flow.
\textbf{15.} Route orders that add to the overnight position by overnight charge as well
as fee, reduce positions at C first, and plan the day's split in advance.
\textbf{16.} Its capacity at B falls to 75 million of position and 300 of settlement:
the plan must move flow and position to C at higher cost, or the fund must trade less.
\textbf{17.} It gives access when no bank will take the fund directly, but adds a layer
of fees, a second set of limits and the intermediary's own credit risk.
\textbf{18.} Its capacity, data use, pre-hedging and last-look practice, window lengths
and whether it trades in the window, in the standard cover sheet; and the same from each
liquidity provider it will face.
\textbf{19.} Named result: \emph{three prime brokers}: flow A 500, B 600, C 100 and
overnight position A 100, B 150, C 50 (USD millions) keep every limit at USD~13\,800 a
day, against 24\,600 with everything at C.
\textbf{20.} A prime broker's credit and a gate that keeps every order inside its limits.
\end{solution}

\begin{solution}{iq:m2:getting-access-fx:1}
It lets the fund trade with many dealers under the prime broker's credit: the fund
executes with a dealer, the trade is given up to the prime broker, which becomes the
dealer's counterparty, and the prime broker books an offsetting trade with the fund,
within limits set in its notice to each dealer.
\iqlookfor{credit intermediation and the three-party structure.}
\end{solution}

\begin{solution}{iq:m2:getting-access-fx:2}
The net open position is the net currency exposure in dollars, here the sum of net long
dollar values across currencies; the \omterm{def:m2:getting-access-fx:limits}{settlement limit} caps the amount to be settled on
each value date, here the sum of the currencies to be received. Compute both from the
book of trades at current rates, per prime broker and per value date.
\iqlookfor{netting by currency and by value date.}
\end{solution}

\begin{solution}{iq:m2:getting-access-fx:3}
In the order path, after the strategy and before the order leaves, with its state kept
in memory and updated on each fill; it must be deterministic and fast, and when it
cannot decide in time, because state is stale or a limit feed is missing, it must
reject: failing closed.
\iqlookfor{placement and fail-closed.}
\end{solution}

\begin{solution}{iq:m2:getting-access-fx:4}
Estimate the flow and position by currency; list what each prime broker allows and its
limits and prices; solve for the cheapest allocation that respects every limit, filling
constrained flows first; keep headroom for stress and for limit cuts; route in real
time with a gate that also considers overnight charges.
\iqlookfor{constraints, costs, headroom.}
\end{solution}

\begin{solution}{iq:m2:getting-access-fx:5}
Its disclosure cover sheet: whether it acts as principal, whether it uses \omterm{def:m2:fx-spot-microstructure:lastlook}{last look} and
how (symmetry, window, trading during the window), whether it pre-hedges, how it uses
client data; then evidence from its own trades with it: fills, rejects, hold times and
mark-outs.
\iqlookfor{disclosures and measured behaviour.}
\end{solution}

\begin{solution}{iq:m2:getting-access-fx:6}
Keep each prime broker's state in fixed arrays indexed by currency and value date,
preallocated; compute an order's effect on the two measures incrementally from its two
legs; partition by prime broker or by strategy to avoid contention, or use one thread
per partition with lock-free queues; no allocation or system calls on the path; measure
tail latency and test against the reference implementation.
\iqlookfor{incremental updates, no allocation, partitioning, tail latency.}
\end{solution}
